%==========================================================================
\section{Feature overview}
In LO-4, every field access and method call must name its member in the source
code. LO-5 reflection lets a program find an object's runtime class, list its
fields and method signatures, read or write a field, call a method, or
construct a new instance of an object, all by a name computed at runtime. Each
reflective access states the type it expects, and the runtime checks it
against the class's metadata before touching the object, aborting on a
mismatch instead of misreading memory. Structural changes, such as defining
classes, adding members or altering dispatch at runtime, are out of scope.
\textbf{Motivating example:} a serializer for any subclass of a common root
class \code{Serial}. It writes out an object's int, bool and String fields by
name, writes \code{Serial}-typed reference fields as object IDs, and rebuilds
each object from its class name with \code{allocate}, so no constructor
convention is needed. Without reflection, every class would need its own save
and load methods and a branch in a factory keyed on class names; with
reflection, a class becomes serializable just by extending \code{Serial}.

\textbf{Code Example:}
\begin{lstlisting}
class Counter (int count;) { int add(int k) { return (count + k); } }
class Main () {
  int main() {
    Counter c; ReflectInfo info;
    c = new Counter(3); info = reflect(c);
    out.print_string(info.class_name());   // "Counter"
    setfield(int, c, "count", 7);
    return invoke(int, c, "add", 5);       // 12
  }
}
\end{lstlisting}

